\documentclass[a4paper,12pt]{article}
\usepackage[T2A]{fontenc}
\usepackage[utf8]{inputenc}
\usepackage[russian]{babel}
\usepackage{amsmath,amssymb,mathtools}
\usepackage{paratype}
\renewcommand{\familydefault}{\sfdefault}
\usepackage{icomma}
\usepackage[left=20mm,right=20mm,top=20mm,bottom=20mm]{geometry}
\usepackage{microtype}
\usepackage{tikz}
\usetikzlibrary{calc,arrows.meta,positioning,shapes.geometric}
\usepackage{tkz-euclide}
\usepackage{tabularx,booktabs,longtable}
\usepackage{enumitem}
\usepackage{hyperref}
\usepackage{parskip}
\usepackage{fancyhdr}
\usepackage{setspace}
\usepackage{xcolor}

\definecolor{accent}{HTML}{287C78}
\definecolor{darkaccent}{HTML}{185451}
\definecolor{textgray}{HTML}{333333}
\definecolor{softgray}{HTML}{F2F4F3}

\hypersetup{colorlinks=true,linkcolor=darkaccent,urlcolor=darkaccent}
\setlength{\parindent}{0pt}
\setlength{\parskip}{0.6em}
\onehalfspacing
\setlist[itemize]{leftmargin=1.5em,itemsep=0.25em,topsep=0.25em}
\setlist[enumerate]{leftmargin=1.8em,itemsep=0.5em,topsep=0.35em}
\setlength{\headheight}{25pt}
\pagestyle{fancy}
\fancyhf{}
\fancyhead[L]{\small\color{textgray}Повторение ЕГЭ}
\fancyhead[R]{\small\color{textgray}07.10.26}
\fancyfoot[C]{\small\color{textgray}\thepage}
\renewcommand{\headrulewidth}{0.3pt}
\renewcommand{\headrule}{\hbox to\headwidth{\color{accent}\leaders\hrule height \headrulewidth\hfill}}

\newcommand{\topic}[1]{%
  \vspace{0.45em}\noindent{\Large\bfseries\color{darkaccent}#1}\par\vspace{0.2em}
  {\color{accent}\rule{\linewidth}{0.6pt}}\vspace{0.35em}
}
\newcommand{\keyidea}[1]{\noindent\textbf{\color{darkaccent}Ключевая идея.} #1}
\newcommand{\warning}[1]{\noindent\textbf{\color{darkaccent}Будьте внимательны.} #1}
\newcommand{\formula}[1]{\[\boxed{#1}\]}

\begin{document}

\begin{titlepage}
\thispagestyle{empty}
\centering
\vspace*{2.5cm}

{\Large\color{accent}ЧЕК-ЛИСТ\par}
\vspace{0.8cm}
{\huge\bfseries\color{darkaccent}Повторение ЕГЭ:\par}
\vspace{0.25cm}
{\LARGE\bfseries\color{darkaccent}быстрые методы в задачах первой части\par}

\vspace{1.4cm}
{\large Алгебраические тождества \textbullet{} тригонометрия \textbullet{} корни\par}
{\large Векторы \textbullet{} классическая вероятность\par}

\vfill
{\large 07.10.26\par}
\vspace{0.7cm}
{\large\bfseries Лёвин Артём Александрович\par}
{\large эксклюзивно для Ксении Клыковой\par}

\vfill
\begin{tikzpicture}[remember picture,overlay]
  \draw[accent!65,line width=1.3pt]
  ([xshift=12mm,yshift=18mm]current page.south west)
  .. controls ([xshift=55mm,yshift=36mm]current page.south west)
  and ([xshift=-55mm,yshift=8mm]current page.south east)
  .. ([xshift=-12mm,yshift=24mm]current page.south east);
\end{tikzpicture}
\end{titlepage}

\topic{1. Не ищите лишние неизвестные: используйте тождества}

\keyidea{Если в условии известны сумма и произведение двух величин, часто можно сразу получить нужную сумму квадратов, не решая систему полностью.}

Основные тождества:
\formula{(x+y)^2=x^2+2xy+y^2}
\formula{(x-y)^2=x^2-2xy+y^2}

Отсюда:
\[
 x^2+y^2=(x+y)^2-2xy,
 \qquad
 x^2+y^2=(x-y)^2+2xy.
\]

\textbf{Пример с прямоугольником.} Пусть стороны равны $x$ и $y$, периметр $P=8$, площадь $S=3{,}5$. Тогда
\[
 x+y=\frac P2=4,\qquad xy=S=3{,}5.
\]
Квадрат диагонали:
\[
 d^2=x^2+y^2=(x+y)^2-2xy=16-7=9,
\]
поэтому $d=3$.

\textbf{Полезная формула для прямоугольника:}
\formula{d^2=\left(\frac P2\right)^2-2S}

\textbf{Ещё один тип.} Если $3k-2t=10$ и $\dfrac{kt}{3}=5$, то $kt=15$. Возводим первое равенство в квадрат:
\[
(3k-2t)^2=9k^2-12kt+4t^2=100.
\]
Значит,
\[
9k^2+4t^2=100+12kt=100+180=280.
\]

\warning{Перед подстановкой отдельных значений $x$ и $y$ спросите себя: действительно ли их нужно искать? Если требуется только комбинация вроде $x^2+y^2$, тождество часто короче и надёжнее.}

\topic{2. Прямоугольный треугольник: используйте дополнительные углы}

В прямоугольном треугольнике $ABC$ с $\angle C=90^\circ$ углы $A$ и $B$ дополняют друг друга до $90^\circ$. Поэтому
\formula{\sin A=\cos B,\qquad \cos A=\sin B}

Если из вершины $C$ опущена высота $CH$ на гипотенузу $AB$, появляются два новых прямоугольных треугольника. Один и тот же угол можно использовать в большом и малом треугольнике.

\begin{center}
\begin{tikzpicture}[scale=0.9]
  \tkzDefPoint(0,0){C}
  \tkzDefPoint(6,0){B}
  \tkzDefPoint(0,4.5){A}
  \tkzDrawPolygon[line width=1pt,color=accent](A,B,C)
  \tkzDefPointBy[projection=onto A--B](C)\tkzGetPoint{H}
  \tkzDrawSegment[line width=1pt,color=accent](C,H)
  \tkzMarkRightAngle[size=0.3](A,C,B)
  \tkzMarkRightAngle[size=0.25](C,H,B)
  \tkzLabelPoints[left](A,C)
  \tkzLabelPoints[right](B)
  \tkzLabelPoints[above right](H)
\end{tikzpicture}
\end{center}

\textbf{Пример.} Пусть $BC=8$ и $\sin A=0{,}5$. Тогда $\cos B=0{,}5$. В треугольнике $BHC$ гипотенуза --- $BC$, поэтому
\[
\cos B=\frac{BH}{BC},\qquad BH=8\cdot0{,}5=4.
\]
Кроме того,
\[
\sin B=\sqrt{1-\cos^2 B}=\frac{\sqrt3}{2},
\qquad
CH=BC\sin B=4\sqrt3.
\]

\warning{Сначала назовите треугольник, в котором используете синус или косинус. Ошибка чаще возникает из-за неверно выбранной гипотенузы или катета.}

\topic{3. Точные корни: проверяйте структуру числа}

Для точного извлечения квадратного корня полезны два приёма.

\textbf{Разложение на множители:}
\[
576=2^6\cdot3^2
\quad\Longrightarrow\quad
\sqrt{576}=2^3\cdot3=24.
\]

\textbf{Десятичную дробь безопаснее переписать как обычную:}
\[
\sqrt{5{,}76}=\sqrt{\frac{576}{100}}=\frac{24}{10}=2{,}4.
\]

Для целого полного квадрата можно быстро оценить диапазон: $20^2<576<30^2$. Последняя цифра $6$ подсказывает кандидатов с единицами $4$ или $6$; проверка даёт $24^2=576$.

\warning{Оценка по последней цифре --- только быстрый фильтр. Финальный результат нужно проверить возведением в квадрат или точным разложением.}

\topic{4. Векторы: координаты, сумма, длина и скалярное произведение}

Для точек $A(x_A,y_A)$ и $B(x_B,y_B)$:
\formula{\overrightarrow{AB}=(x_B-x_A;\ y_B-y_A)}

Сумма векторов считается покоординатно:
\formula{(x_1;y_1)+(x_2;y_2)=(x_1+x_2;\ y_1+y_2)}

Длина и квадрат длины:
\formula{|\vec a|=\sqrt{x_a^2+y_a^2},\qquad |\vec a|^2=x_a^2+y_a^2}

Скалярное произведение:
\formula{\vec a\cdot\vec b=x_ax_b+y_ay_b}

\textbf{Пример.} Пусть
\[
\vec a=(-10;3),\qquad \vec b=(-1;-6),\qquad \vec c=(-2;6).
\]
Тогда
\[
\vec a+\vec b=(-11;-3),
\]
и
\[
(\vec a+\vec b)\cdot\vec c=(-11)(-2)+(-3)\cdot6=22-18=4.
\]

Если $\vec a=(-3;1)$ и $\vec b=(3;2)$, то
\[
\vec a+\vec b=(0;3),\qquad |\vec a+\vec b|^2=0^2+3^2=9.
\]

\warning{Нельзя складывать длины векторов вместо самих векторов. Сначала найдите координаты суммы, затем работайте с её длиной.}

\topic{5. Классическая вероятность: зафиксируйте одного участника}

Если все исходы равновозможны,
\formula{P=\frac{\text{число благоприятных исходов}}{\text{общее число исходов}}}

\textbf{Группы.} $16$ человек случайно делят на группы по $4$. Зафиксируем Анну. Для Татьяны остаётся $15$ возможных мест, из них $3$ --- в группе Анны. Поэтому
\[
P=\frac{3}{15}=\frac15.
\]

\textbf{Соперник.} В турнире $76$ шахматистов, из них $19$ из России, включая Ивана. У Ивана $75$ возможных соперников, из них $18$ --- российские шахматисты. Поэтому
\[
P=\frac{18}{75}=\frac{6}{25}.
\]

\warning{После фиксации одного человека не считайте его снова среди возможных партнёров или соперников. Именно здесь часто появляется ошибка на единицу.}

\topic{6. Что держать в памяти перед первой частью ЕГЭ}

\begin{itemize}
  \item Если известны $x+y$ и $xy$, ищите нужное выражение через тождество.
  \item В прямоугольном треугольнике используйте $\sin A=\cos B$ и $\cos A=\sin B$.
  \item Перед вычислением корня оцените порядок числа и проверьте, является ли оно полным квадратом.
  \item Координаты вектора --- это смещение; сумма векторов считается покоординатно.
  \item Для $|\vec a|^2$ корень извлекать не нужно.
  \item В вероятности сначала чётко назовите благоприятные и все равновозможные исходы.
  \item В короткой задаче ищите короткий путь, но каждый переход должен быть математически обоснован.
\end{itemize}

\clearpage
\topic{Тренировочный блок --- Путь А}

\textbf{10 заданий. Решения в этом блоке не записаны. Сложность постепенно возрастает.}

\begin{enumerate}
  \item Периметр прямоугольника равен $14$, а площадь равна $12$. Найдите диагональ прямоугольника.

  \item Известно, что $2a+3b=11$ и $ab=2$. Найдите $4a^2+9b^2$.

  \item Известно, что $4m-n=7$ и $mn=3$. Найдите $16m^2+n^2$.

  \item В прямоугольном треугольнике $ABC$ угол $C$ равен $90^\circ$, $CH$ --- высота к гипотенузе $AB$. Известно, что $AC=10$ и $\cos A=0{,}6$. Найдите $AH$.

  \item В прямоугольном треугольнике $ABC$ угол $C$ равен $90^\circ$, $CH$ --- высота к гипотенузе $AB$. Известно, что $BC=10$ и $\sin A=0{,}6$. Найдите $CH$.

  \item Вычислите $\sqrt{2304}$ без калькулятора.

  \item Даны векторы $\vec a=(-4;5)$, $\vec b=(1;-2)$ и $\vec c=(3;4)$. Найдите $(\vec a+\vec b)\cdot\vec c$.

  \item Даны точки $A(-2;5)$ и $B(4;1)$. Найдите $|\overrightarrow{AB}|^2$.

  \item $24$ ученика случайно распределяются по группам по $6$ человек. Какова вероятность, что два заранее выбранных ученика окажутся в одной группе?

  \item В турнире участвуют $120$ человек, из них $32$ представляют один регион. Один из этих $32$ участников выбирает себе случайного соперника из остальных. Найдите вероятность, что соперник будет из того же региона.
\end{enumerate}

\clearpage
\topic{Ответы}

\renewcommand{\arraystretch}{1.35}
\begin{center}
\begin{tabularx}{0.88\linewidth}{@{}>{\bfseries}c X@{}}
\toprule
№ & Ответ \\
\midrule
1 & $5$ \\
2 & $97$ \\
3 & $73$ \\
4 & $6$ \\
5 & $8$ \\
6 & $48$ \\
7 & $3$ \\
8 & $52$ \\
9 & $\dfrac{5}{23}$ \\
10 & $\dfrac{31}{119}$ \\
\bottomrule
\end{tabularx}
\end{center}

\vfill
\begin{center}
\small\color{textgray}
Лёвин Артём Александрович эксклюзивно для Ксении Клыковой
\end{center}

\end{document}